\documentclass[11pt]{article}
\usepackage[margin=1in]{geometry}
\usepackage{times}
\usepackage{hyperref}
\hypersetup{colorlinks=true, linkcolor=black, urlcolor=blue, citecolor=black}
\title{A Feasibility Study Is Not a Discovery: Boeing's GRASP Program and the Antigravity Rumor It Generated}
\author{Flyxion \\ \small Independent Researcher}
\date{}
\begin{document}
\maketitle

\begin{abstract}
Reports in the early 2000s that Boeing's Phantom Works division had established a research effort, popularly referred to as GRASP, to investigate gravity modification circulated widely as evidence that a major aerospace contractor had found a credible path to gravity control propulsion. This paper argues that the available public record supports a considerably more modest reading: an internal feasibility review of an already speculative and largely discredited literature, undertaken as routine technology-scouting due diligence, which found nothing worth pursuing and was quietly discontinued, and that the persistence of the antigravity reading owes to the gap between what a large company's internal exploratory review actually is and what public reporting on leaked or rumored internal programs tends to imply.
\end{abstract}

\section{Introduction}

Trade press reporting around 2002, most notably in Jane's Defence Weekly, described a Boeing research effort referred to internally as GRASP, said to be investigating gravity modification technologies including some of the claims surveyed elsewhere in this series, such as Podkletnov's rotating superconductor work. The reporting, and much of the subsequent circulation of it, framed the existence of this internal effort as evidence that a serious aerospace engineering organization considered gravity control technically credible, and by implication as third-party validation of the underlying antigravity claims themselves.

\section{What Large Engineering Organizations Actually Do With Speculative Claims}

Aerospace and defense contractors of Boeing's size maintain technology-scouting functions whose explicit purpose is to review emerging and even speculative claims across a wide range of fields, precisely so that a genuine breakthrough is not missed and so that competitors cannot claim exclusive awareness of a promising lead. The overwhelming majority of what such a function reviews does not pan out, and a review being conducted is evidence only that the claim was judged worth a bounded, low-cost look, not that the claim was judged likely to be true. Public reporting on internal corporate research efforts, especially when sourced from limited leaked documentation or from statements by researchers eager for their work to be taken seriously, systematically overstates the significance of a company having looked at something relative to a company having found something.

\section{What Actually Happened, So Far as the Public Record Shows}

No published Boeing patent, technical paper, or flight demonstration has ever resulted from the GRASP effort, and subsequent reporting, including follow-up coverage after the initial trade press story, indicated the program was discontinued without producing findings supportive of the underlying gravity modification claims it had been formed to evaluate. This is the expected outcome of a due-diligence review of a literature that, as the other essays in this series argue in detail, does not identify a physical mechanism capable of doing what it claims, and a review reaching this conclusion and quietly winding down is not a newsworthy non-event in the way that a review reaching a positive conclusion would be, which likely accounts for the asymmetry between how widely the program's existence was reported and how little attention its quiet discontinuation received.

\section{The Entropic Gravity Framing}

Nothing in the GRASP program's reported scope, insofar as the public record describes it, proposed a novel mechanism for gravitational coupling beyond the ones already addressed in this series, primarily the Podkletnov superconductor claim. A corporate feasibility review evaluating a mechanism that has no available channel to affect gravitational curvature, on either the standard general-relativistic or the entropic account, would be expected to conclude exactly what the available record suggests it concluded, regardless of how the review was initially reported.

\section{Objections}

It might be argued that a company would have strong incentive to keep a genuinely successful gravity modification result confidential for competitive and national security reasons, and that the absence of published results is therefore uninformative. This is a general-purpose objection available to defend any claim of suppressed success, and it proves too much: it would equally excuse the total absence of results for every device surveyed in this series, converting a complete lack of evidence into an argument for secrecy rather than for absence, a move that abandons ordinary evidentiary standards rather than satisfying them.

\section{Conclusion}

The public record surrounding Boeing's GRASP program supports a mundane, if slightly deflating, conclusion: a large engineering organization performed routine due diligence on an already weak literature, found nothing to pursue, and discontinued the effort, a sequence of events considerably less remarkable than its circulation as an antigravity validation story has suggested.

\end{document}
